\begin{figure}[!ht]
  \centering
  \begin{tikzpicture}[
    node distance=5mm and 7mm,
    comparisonstep/.style={
      draw=black!65,
      rounded corners=1.5pt,
      line width=.45pt,
      text width=5.7cm,
      align=center,
      inner xsep=5pt,
      inner ysep=4pt,
      font=\small
    },
    comparisonwide/.style={comparisonstep,text width=12.45cm},
    dependency/.style={-{Stealth[length=2mm,width=1.4mm]},line width=.55pt}
  ]
    \node[comparisonstep,minimum height=12mm] (upperfree) {
      \textbf{Upper terminal kernels are free}\\
      at every $y>u$, in both systems.
    };
    \node[comparisonstep,minimum height=12mm,right=of upperfree] (upperequal) {
      \textbf{Upper polynomials agree}\\
      $P_{y,b}=P'_{\iota(y),b'}$ for every $y>u$.
    };
    \node[comparisonstep,minimum height=16mm,below=of upperfree] (comparison) {
      \textbf{Hilbert comparison}\\
      The inclusion $I\subseteq f_eJ$ bounds each increase in Hilbert series.
    };
    \node[comparisonstep,minimum height=16mm,below=of upperequal] (reciprocity) {
      \textbf{Reciprocity and cancellation}\\
      Cancel the upper terms; compare in both directions.
    };
    \path (comparison.south west) -- (reciprocity.south east)
      coordinate[midway] (middle);
    \node[comparisonwide,below=6mm of middle] (polynomial) {
      \textbf{Polynomial equality at the bottom}\\
      The strict degree bound gives $P_{u,b}=P'_{u',b'}$.
    };
    \node[comparisonwide,below=of polynomial] (images) {
      \textbf{Equality of all edge images}\\
      Exact endpoint comparison forces every bound to be attained: $I=f_eJ$.
    };
    \node[comparisonwide,below=of images] (freeness) {
      \textbf{Bottom terminal kernels are free}\\
      Reverse the bottom schedule, starting from $L_u(\emptyset)=B^u$.
    };
    \draw[dependency] (upperfree) -- (comparison);
    \draw[dependency] (upperequal) -- (reciprocity);
    \draw[dependency] (comparison) -- (reciprocity);
    \draw[dependency] (reciprocity.south) -- (reciprocity.south |- polynomial.north);
    \draw[dependency] (polynomial) -- (images);
    \draw[dependency] (images) -- (freeness);
  \end{tikzpicture}
  \caption{The two uses of smaller-rank induction. Upper freeness supplies
    the edge bound; upper polynomial equality permits row cancellation
    against reciprocity. Polynomial equality then forces the edge bounds
    to be exact, and the reverse bottom schedule proves terminal-kernel
    freeness.}
  \label{fig:comparison-dependencies}
\end{figure}
